% FILE: chapters/ch01_s1_python_fundamentals.tex
\chapter{Programming and Tools}\label{ch:programming_tools}
\section{Python Fundamentals}\label{sec:python_fundamentals}

Machine learning is not magic. At its heart, it is mathematics turned into working software. Before an algorithm can predict the temperature of a greenhouse, detect a failing motor bearing, or identify a leaf disease on an agricultural drone, we need a language to express our thoughts to the computer. 

Python is the universal language of modern artificial intelligence and data science. It reads almost like plain English, allowing us to focus on understanding concepts rather than fighting confusing syntax. In this section, we build our programming foundation from absolute zero. If you have never written a line of code in your life, you are in the right place.

\begin{figure}[htbp]
\centering
\begin{tikzpicture}[
    node distance=1.6cm,
    box/.style={rectangle, draw=keyideablue, fill=keyideablue!10, rounded corners=2mm, thick, text width=2.4cm, align=center, minimum height=1.1cm, font=\small\sffamily},
    arrow/.style={-{Stealth[length=2.5mm]}, thick, draw=keyideablue}
]
    \node[box] (sensor) {Physical Sensor\\(e.g., DHT22, BME280)};
    \node[box, right=of sensor] (mcu) {Microcontroller\\(ESP32 / ESP8266)};
    \node[box, right=of mcu] (python) {Python Data Pipeline\\(Clean \& Structure)};
    \node[box, right=of python] (ml) {Machine Learning\\(Predict \& Act)};

    \draw[arrow] (sensor) -- node[above, font=\tiny\sffamily] {Voltage / Bits} (mcu);
    \draw[arrow] (mcu) -- node[above, font=\tiny\sffamily] {Serial / Wi-Fi} (python);
    \draw[arrow] (python) -- node[above, font=\tiny\sffamily] {Tensors} (ml);
\end{tikzpicture}
\caption{The end-to-end data flow: from physical measurement to Python processing and machine learning.}\label{fig:iot_pipeline}
\end{figure}

% ==============================================================================
% 1. VARIABLES, TYPES, AND ARITHMETIC
% ==============================================================================
\subsection{Variables, Types, and Basic Arithmetic}

\begin{keyidea}[Variables as Labelled Boxes]
A \textbf{variable} is a named storage location in your computer's memory that holds a piece of data. Instead of remembering the raw memory address (like \texttt{0x7ffee3}), we give it a clear, human-readable name, such as \texttt{temperature} or \texttt{node\_id}.
\end{keyidea}

\paragraph{Why It Exists: The Real-World Problem}
Physical sensors continuously measure the environment: temperature, humidity, pressure, or electrical current. To calculate averages, check for anomalies, or train a model, a computer must temporarily hold these readings in memory. Variables give us a way to store, update, and reuse values throughout our program.

\begin{analogy}[Pantry Storage Containers]
Think of your computer's memory as a kitchen pantry. If you pour sugar, flour, and coffee into unmarked bowls, you will quickly ruin your recipe. Instead, you stick a label on each jar: ``Flour'', ``White Sugar'', or ``Dark Roast''. In Python, the label on the jar is the variable name, and whatever you pour inside is the value.
\end{analogy}

\paragraph{The Idea in Words, Step by Step}
In Python, you create a variable using the assignment operator, written as the equals sign (\texttt{=}).
\begin{enumerate}
    \item You write a variable name on the left side (for example, \texttt{temperature}).
    \item You place the assignment operator \texttt{=} in the middle.
    \item You write the value on the right side (for example, \texttt{23.5}).
\end{enumerate}
Python automatically inspects the value and determines its \textbf{data type}:
\begin{itemize}
    \item \textbf{Integer (\code{int})}: Whole numbers without a decimal point, such as \texttt{42}, \texttt{0}, or \texttt{-12}. Used for counting items, device IDs, or epoch iterations.
    \item \textbf{Floating-point (\code{float})}: Numbers with a fractional decimal part, such as \texttt{3.14159}, \texttt{23.5}, or \texttt{0.001}. Used for continuous measurements like voltage, loss, and weights.
    \item \textbf{String (\code{str})}: Plain text enclosed in single or double quotes, such as \texttt{"sensor\_01"} or \texttt{"Normal"}.
    \item \textbf{Boolean (\code{bool})}: Represents truth values: exactly \texttt{True} or \texttt{False} (capitalized in Python). Used for flags, alarms, and binary decisions.
\end{itemize}

\paragraph{The Maths: Sensor Voltage Conversion}
Microcontrollers measure physical quantities by converting an analog voltage into a digital number using an Analog-to-Digital Converter (ADC). On an ESP32, the ADC reads 12 bits of precision, producing an integer between $0$ and $4095$ representing a voltage between $0\text{ V}$ and $3.3\text{ V}$.

The formula to convert an integer ADC reading ($ADC$) to a physical voltage ($V$) is:
\begin{equation}
V = \frac{ADC}{ADC_{\max}} \times V_{\text{ref}} = \frac{ADC}{4095} \times 3.3\label{eq:adc_conversion}
\end{equation}
Let us calculate this by hand for a raw reading of $ADC = 2048$:
\begin{enumerate}
    \item Divide the raw integer by the maximum resolution: $\frac{2048}{4095} \approx 0.500122$.
    \item Multiply by reference voltage $3.3\text{ V}$: $0.500122 \times 3.3 \approx 1.6504\text{ V}$.
\end{enumerate}
Our floating-point variable will hold the value $1.6504$.

\paragraph{Python Code}
\begin{lstlisting}[language=Python]
# Defining sensor parameters
sensor_name = "DHT22_Greenhouse"  # str: name of the sensor
raw_adc = 2048                    # int: 12-bit ADC reading (0 to 4095)
v_ref = 3.3                       # float: operating reference voltage in Volts
is_active = True                  # bool: sensor status flag

# Calculating voltage using Equation 1.1
voltage = (raw_adc / 4095.0) * v_ref

# Formatted output using f-strings
print(f"Sensor: {sensor_name}")
print(f"Raw Reading: {raw_adc} (Type: {type(raw_adc).__name__})")
print(f"Calculated Voltage: {voltage:.4f} V (Type: {type(voltage).__name__})")
print(f"Node Online: {is_active}")
\end{lstlisting}

\begin{lstlisting}[style=outputstyle]
Sensor: DHT22_Greenhouse
Raw Reading: 2048 (Type: int)
Calculated Voltage: 1.6504 V (Type: float)
Node Online: True
\end{lstlisting}

\paragraph{Line-by-Line Explanation}
\begin{itemize}
    \item Lines 2--5: We create four variables. Python infers the type of each variable automatically at runtime.
    \item Line 8: We perform division and multiplication. Because \code{4095.0} and \code{v\_ref} are floats, Python promotes the result of \code{voltage} to a \code{float}.
    \item Lines 11--14: We use an \textbf{f-string} (formatted string literal, starting with \code{f"..."}). Placing curly braces \code{\{voltage:.4f\}} formats the float to exactly four decimal places.
\end{itemize}

\begin{iotbox}[Raw ADC Readings on ESP32 Microcontrollers]
When programming an ESP32 in MicroPython or receiving serial streams in C++, you always begin with raw ADC values. For example, pin GPIO34 connected to a soil moisture sensor might stream numbers between $1200$ (sopping wet soil) and $3800$ (bone dry soil). Python scripts running on your laptop or Raspberry Pi gateway immediately cast these readings into physical percentages.
\end{iotbox}

\begin{warning}[String Concatenation vs. Addition]
A common mistake when reading data from serial ports or text files is forgetting to convert strings into numbers:
\begin{lstlisting}[language=Python]
reading_a = "25"  # This is a string (text), not a number
reading_b = "5"
print(reading_a + reading_b)  # Outputs "255", NOT 30!
\end{lstlisting}
The \code{+} operator glues two strings together (called \textbf{concatenation}). To perform mathematical addition, explicitly convert both operands with \code{int()} or \code{float()}: \code{float(reading_a) + float(reading_b)}.
\end{warning}

\begin{takeaways}
\begin{itemize}
    \item Variables store data in memory and are created using the assignment operator (\code{=}).
    \item The four primary primitive types in Python are \code{int}, \code{float}, \code{str}, and \code{bool}.
    \item Data arriving from sensors or text files is almost always received as strings and must be cast using \code{int()} or \code{float()} before doing math.
\end{itemize}
\end{takeaways}

% ==============================================================================
% 2. CONDITIONALS AND DECISION MAKING
% ==============================================================================
\subsection{Conditionals and Decision Making}

\begin{keyidea}[Conditional Branching]
A \textbf{conditional statement} allows a program to inspect a condition and choose which path of code to execute. If a condition evaluates to \code{True}, one block of code runs; if \code{False}, an alternative block can run instead.
\end{keyidea}

\paragraph{Why It Exists: The Real-World Problem}
Static programs that perform the exact same steps every second are useless in dynamic environments. If a greenhouse reaches $38^\circ\text{C}$, the ventilation fans must turn on; if the temperature drops to $2^\circ\text{C}$, the heating lamps must activate. Automated systems must make decisions dynamically without human intervention.

\begin{analogy}[A Highway Fork]
Imagine driving a car toward a highway fork. If your dashboard navigation shows your exit is on the right, you steer right. Otherwise, you stay in the main lane. You do not drive on both roads at the same time; the state of your journey determines which single road you travel down.
\end{analogy}

\paragraph{The Idea in Words, Step by Step}
Conditional statements in Python rely on three keywords: \code{if}, \code{elif} (short for ``else if''), and \code{else}.
\begin{enumerate}
    \item The program evaluates the boolean condition immediately following the \code{if}.
    \item If the expression is \code{True}, the indented block of code underneath runs, and Python skips all remaining branches.
    \item If the condition is \code{False}, Python steps to the next \code{elif} statement and checks its condition.
    \item If none of the conditions evaluate to \code{True}, the fallback \code{else} block runs.
\end{enumerate}

Python uses comparison operators:
\begin{itemize}
    \item Equal to: \code{==} (two equals signs; never use single \code{=}, which assigns values).
    \item Not equal to: \code{!=}.
    \item Greater than / Less than: \code{>}, \code{<}, \code{>=}, \code{<=}.
\end{itemize}
Multiple conditions can be combined using logical operators: \code{and} (both must be true), \code{or} (at least one must be true), and \code{not} (reverses the boolean value).

\paragraph{Python Code}
\begin{lstlisting}[language=Python]
temperature = 34.2  # Degrees Celsius
humidity = 82.0     # Percentage (%)

# Environmental monitoring logic
if temperature > 35.0:
    status = "CRITICAL: Extreme Heat Warning!"
    exhaust_fan = True
elif temperature > 30.0 and humidity > 75.0:
    status = "WARNING: High Heat and Humidity - Mold Risk"
    exhaust_fan = True
elif temperature < 10.0:
    status = "ALERT: Cold Temperature - Frost Risk"
    exhaust_fan = False
else:
    status = "NOMINAL: Conditions Optimal"
    exhaust_fan = False

print(f"Status Report: {status}")
print(f"Exhaust Fan Active: {exhaust_fan}")
\end{lstlisting}

\begin{lstlisting}[style=outputstyle]
Status Report: WARNING: High Heat and Humidity - Mold Risk
Exhaust Fan Active: True
\end{lstlisting}

\paragraph{Line-by-Line Explanation}
\begin{itemize}
    \item Lines 1--2: We initialize temperature and humidity values.
    \item Line 5: Python tests \code{34.2 > 35.0}. This evaluates to \code{False}, so the first block is skipped.
    \item Line 8: Python tests \code{34.2 > 30.0 and 82.0 > 75.0}. Both sub-conditions are \code{True}, so the combined expression is \code{True}.
    \item Lines 9--10: The warning message is assigned, and the fan flag is set to \code{True}. Python now skips lines 11--16 completely.
\end{itemize}

\begin{iotbox}[Edge Actuation Logic]
In IoT, conditionals directly drive hardware actuators. When an ESP8266 or ESP32 running MicroPython detects an alert threshold, setting \code{exhaust\_fan = True} corresponds to writing a high voltage to a relay pin: \code{relay\_pin.value(1)}. This closes an electrical circuit, powering on a physical $240\text{ V}$ industrial ventilation fan.
\end{iotbox}

\begin{warning}[Indentation Defines Structure]
Unlike C, C++, or Java, which use curly braces \texttt{\{ \}} to enclose blocks of code, Python uses \textbf{indentation} (spaces). Standard Python uses four spaces per indentation level. Inconsistent tabs and spaces will produce an \code{IndentationError}.
\end{warning}

\begin{takeaways}
\begin{itemize}
    \item Branching statements allow programs to react dynamically to incoming sensor data.
    \item Use \code{==} to compare two values for equality; never use \code{=}.
    \item Indentation is not decorative in Python; it defines the scope and hierarchy of code blocks.
\end{itemize}
\end{takeaways}

% ==============================================================================
% 3. LOOPS: ITERATING OVER DATA
% ==============================================================================
\subsection{Loops: Iterating Over Data Streams}

\begin{keyidea}[Repetition via Loops]
A \textbf{loop} executes a block of code multiple times. A \code{for} loop iterates over a fixed sequence of items, while a \code{while} loop repeats as long as a specified condition remains \code{True}.
\end{keyidea}

\paragraph{Why It Exists: The Real-World Problem}
In machine learning and IoT, data never arrives as a single value. It arrives in arrays, lists, time-series streams, and tables with thousands of rows. Writing out separate lines of code for each measurement is impossible. Loops allow us to apply the exact same mathematical transformation across millions of data points effortlessly.

\begin{analogy}[The Conveyor Belt Inspector]
Imagine standing beside a conveyor belt in an agricultural packing plant. Apples roll past you one by one. For every apple, you inspect its diameter: if it is too small, you sort it into juice; if it is large, you pack it for export. You perform the exact same test repeatedly until the belt stops. That is a \code{for} loop.
\end{analogy}

\paragraph{The Idea in Words, Step by Step}
Python offers two main loops:
\begin{enumerate}
    \item \textbf{\code{for} loop}: Reads items sequentially from an iterable (a collection like a list or a range of numbers) until the sequence ends.
    \item \textbf{\code{while} loop}: Continues repeating until its conditional expression turns \code{False}.
\end{enumerate}
To control loops, Python provides:
\begin{itemize}
    \item \code{break}: Immediately halts the loop and jumps to the code following it.
    \item \code{continue}: Skips the rest of the current iteration and jumps directly to the next item.
    \item \code{enumerate()}: Supplies both the current loop index ($0, 1, 2, \dots$) and the item itself.
    \item \code{zip()}: Steps through two or more collections simultaneously in parallel pairs.
\end{itemize}

\paragraph{The Maths: Calculating the Sample Mean}
The sample mean (arithmetic average), denoted by $\bar{x}$, is the foundation of machine learning statistics. For a set of $N$ sensor observations $x_1, x_2, \dots, x_N$, the formula is:
\begin{equation}
\bar{x} = \frac{1}{N}\sum_{i=1}^N x_i = \frac{x_1 + x_2 + \dots + x_N}{N}\label{eq:sample_mean}
\end{equation}
Let us calculate this by hand for $N=4$ temperature measurements: $22.0, 24.0, 23.0, 27.0$.
\begin{enumerate}
    \item Sum the values: $\sum = 22.0 + 24.0 + 23.0 + 27.0 = 96.0$.
    \item Divide by count $N=4$: $\bar{x} = \frac{96.0}{4} = 24.0^\circ\text{C}$.
\end{enumerate}

\paragraph{Python Code}
\begin{lstlisting}[language=Python]
# A sequence of hourly soil moisture sensor readings (%)
readings = [45.2, 44.8, -1.0, 46.1, 45.9]  # -1.0 is a known sensor error code

total_moisture = 0.0
valid_count = 0

print("Processing incoming telemetry stream:")
for index, reading in enumerate(readings):
    # Skip corrupted sensor values using continue
    if reading < 0.0:
        print(f"  [Index {index}] Malformed packet ({reading}%). Discarding.")
        continue
    
    total_moisture += reading
    valid_count += 1
    print(f"  [Index {index}] Valid: {reading}% (Running sum: {total_moisture:.1f})")

# Calculate arithmetic mean using Equation 1.2
if valid_count > 0:
    mean_moisture = total_moisture / valid_count
    print(f"\nFinal Valid Samples: {valid_count}")
    print(f"Average Soil Moisture: {mean_moisture:.2f}%")
\end{lstlisting}

\begin{lstlisting}[style=outputstyle]
Processing incoming telemetry stream:
  [Index 0] Valid: 45.2% (Running sum: 45.2)
  [Index 1] Valid: 44.8% (Running sum: 90.0)
  [Index 2] Malformed packet (-1.0%). Discarding.
  [Index 3] Valid: 46.1% (Running sum: 136.1)
  [Index 4] Valid: 45.9% (Running sum: 182.0)

Final Valid Samples: 4
Average Soil Moisture: 45.50%
\end{lstlisting}

\paragraph{Line-by-Line Explanation}
\begin{itemize}
    \item Line 2: A Python list holds five numbers, including an invalid $-1.0$ reading common when a wire disconnects.
    \item Line 8: \code{enumerate(readings)} pairs every element with its integer position: $(0, 45.2), (1, 44.8), \dots$
    \item Line 10: The conditional checks if the reading is negative.
    \item Line 12: When \code{reading == -1.0}, \code{continue} triggers. Lines 14--16 are skipped for that item, advancing directly to index 3.
    \item Lines 19--21: We verify \code{valid\_count > 0} to prevent dividing by zero, then calculate $\bar{x}$.
\end{itemize}

\begin{warning}[Infinite Loops]
A \code{while} loop whose condition never changes will freeze your program and consume $100\%$ of a CPU core:
\begin{lstlisting}[language=Python]
count = 0
while count < 5:
    print(count)
    # If you forget "count += 1", this loop will run forever!
\end{lstlisting}
Always verify that your loop's state changes on every single iteration toward its stopping condition.
\end{warning}

\begin{takeaways}
\begin{itemize}
    \item Use \code{for} loops when iterating over known collections or ranges.
    \item Use \code{while} loops when repeating until an unpredictable external condition occurs (e.g., listening for a network signal).
    \item \code{enumerate()} gives simultaneous access to array indexes and values.
\end{itemize}
\end{takeaways}

% ==============================================================================
% 4. FUNCTIONS: MODULAR AND REUSABLE LOGIC
% ==============================================================================
\subsection{Functions: Modular and Reusable Logic}

\begin{keyidea}[Functions as Reusable Machines]
A \textbf{function} is a self-contained block of code that accepts inputs (called \textbf{arguments}), performs a specific set of operations, and returns a result. It allows us to write logic once and reuse it everywhere.
\end{keyidea}

\paragraph{Why It Exists: The Real-World Problem}
In machine learning, you must repeatedly scale features, clean text, compute errors, and evaluate metrics across training sets, test sets, and live sensor feeds. If you copy and paste twelve lines of math across your codebase, fixing a tiny bug later means locating and editing every copy by hand. Functions make software modular, readable, and easy to maintain.

\begin{analogy}[The Automatic Bread Maker]
Think of a function as a tabletop bread machine. You pour in flour, yeast, and water through the top hatch (the input parameters). Inside, a pre-programmed sequence kneads, proves, and bakes the dough. An hour later, out pops a fresh loaf of bread (the return value). You do not need to rebuild the machine every time you want bread; you just supply ingredients and take the output.
\end{analogy}

\paragraph{The Idea in Words, Step by Step}
\begin{enumerate}
    \item A function is declared using the \code{def} keyword, followed by the function name, parentheses containing parameters, and a colon: \code{def compute\_area(width, height):}.
    \item An optional, highly recommended \textbf{docstring} (triple quotes \texttt{"""}) explains what the function does.
    \item The indented body contains the processing steps.
    \item The \code{return} statement passes the final computed value back to whatever part of the program called the function. If omitted, Python returns \code{None}.
    \item Parameters can have \textbf{default values}, making them optional when called.
\end{enumerate}

\paragraph{The Maths: Min-Max Feature Scaling}
Machine learning algorithms struggle when inputs have vastly different numeric scales (for example, temperature between $10$ and $40^\circ\text{C}$ paired with atmospheric pressure between $95{,}000$ and $105{,}000\text{ Pa}$). \textbf{Min-Max Scaling} compresses any numeric range into normalized values strictly between $0.0$ and $1.0$:
\begin{equation}
x_{\text{scaled}} = \frac{x - x_{\min}}{x_{\max} - x_{\min}}\label{eq:min_max_scaling}
\end{equation}
Let us normalize a live temperature reading $x = 25.0^\circ\text{C}$ given an expected sensor minimum $x_{\min} = 10.0^\circ\text{C}$ and maximum $x_{\max} = 30.0^\circ\text{C}$:
\begin{enumerate}
    \item Subtract the minimum: $x - x_{\min} = 25.0 - 10.0 = 15.0$.
    \item Calculate the total range: $x_{\max} - x_{\min} = 30.0 - 10.0 = 20.0$.
    \item Divide: $x_{\text{scaled}} = \frac{15.0}{20.0} = 0.75$.
\end{enumerate}
The normalized value is $0.75$.

\paragraph{Python Code}
\begin{lstlisting}[language=Python]
def min_max_scale(value: float, min_val: float = 0.0, max_val: float = 100.0) -> float:
    """
    Normalizes a numerical reading into the range [0.0, 1.0].
    
    Formula: (value - min_val) / (max_val - min_val)
    """
    if min_val >= max_val:
        raise ValueError(f"min_val ({min_val}) must be strictly less than max_val ({max_val})")
    
    # Calculate normalization using Equation 1.3
    scaled = (value - min_val) / (max_val - min_val)
    return scaled

# Using the function with positional and keyword arguments
temp_scaled = min_max_scale(25.0, min_val=10.0, max_val=30.0)
humidity_scaled = min_max_scale(80.0)  # Uses defaults: min_val=0.0, max_val=100.0

print(f"Scaled Temperature (25 C in [10, 30]): {temp_scaled:.4f}")
print(f"Scaled Humidity    (80 % in [0, 100]): {humidity_scaled:.4f}")
\end{lstlisting}

\begin{lstlisting}[style=outputstyle]
Scaled Temperature (25 C in [10, 30]): 0.7500
Scaled Humidity    (80 % in [0, 100]): 0.8000
\end{lstlisting}

\paragraph{Line-by-Line Explanation}
\begin{itemize}
    \item Line 1: We define \code{min\_max\_scale}. We add optional \textbf{type hints} (\code{value: float}) and specify default values (\code{min\_val = 0.0}).
    \item Lines 2--5: The docstring documents the function's contract and mathematical basis.
    \item Lines 6--7: Guard clause. If someone provides nonsensical thresholds where $x_{\min} \ge x_{\max}$, we stop before dividing by zero.
    \item Line 10: We return the computed floating-point value.
    \item Line 15: Because we omit the optional arguments for humidity, Python automatically uses $0.0$ and $100.0$.
\end{itemize}

\begin{iotbox}[Edge Normalization]
On microcontrollers running TinyML neural networks, inputs must be normalized before feeding them to the model tensors. Writing a compact scaling function ensures that the exact same math applied during Python model training on a cloud server runs on the ESP32 in the field.
\end{iotbox}

\begin{takeaways}
\begin{itemize}
    \item Functions encapsulate logic, eliminate duplication, and enhance testability.
    \item The \code{return} statement sends values back; without it, functions evaluate to \code{None}.
    \item Default parameter values allow flexible calling styles while preserving safe fallbacks.
\end{itemize}
\end{takeaways}

% ==============================================================================
% 5. DATA STRUCTURES: LISTS, TUPLES, DICTIONARIES, AND SETS
% ==============================================================================
\subsection{Data Structures: Lists, Tuples, Dictionaries, and Sets}

\begin{keyidea}[Built-in Data Collections]
Python provides four foundational collections for organizing multiple pieces of data in memory:
\begin{itemize}
    \item \textbf{List}: An ordered, mutable (modifiable) sequence of elements.
    \item \textbf{Tuple}: An ordered, immutable (read-only) sequence of elements.
    \item \textbf{Dictionary}: An unordered (insertion-ordered in modern Python) map of unique keys to values.
    \item \textbf{Set}: An unordered collection of unique, non-duplicate items.
\end{itemize}
\end{keyidea}

\paragraph{Why It Exists: The Real-World Problem}
A sensor does not simply output isolated numbers; it outputs a record containing an identification tag, a timestamp, a geographical coordinate, and an array of channel readings. We need structured ways to store and look up related data points quickly.

\begin{analogy}[Everyday Organizing Tools]
\begin{itemize}
    \item \textbf{List}: A shopping list written on scrap paper. You can add items, cross items off, or reorder them.
    \item \textbf{Tuple}: The GPS coordinates stamped onto a stone monument: $(45.5017^\circ\text{N}, 73.5673^\circ\text{W})$. You cannot change the stone's coordinates without destroying it.
    \item \textbf{Dictionary}: A pocket language dictionary or smartphone contacts list. You search for a unique name (the key) to immediately find their phone number (the value).
    \item \textbf{Set}: A bag of colored marbles. If you drop in a blue marble, and a blue marble is already inside, the bag still contains only one blue marble.
\end{itemize}
\end{analogy}

\begin{table}[htbp]
\centering
\caption{Comparison of Core Python Data Structures}\label{tab:python_collections}
\begin{tabular}{@{}lllll@{}}
\toprule
\textbf{Structure} & \textbf{Syntax} & \textbf{Ordered?} & \textbf{Mutable?} & \textbf{Duplicates?} \\ \midrule
\textbf{List}       & \code{[1, 2, 3]}      & Yes               & Yes               & Yes                   \\
\textbf{Tuple}      & \code{(1, 2, 3)}      & Yes               & No                & Yes                   \\
\textbf{Dictionary} & \code{\{"k": "v"\}}   & Yes (by insert)   & Yes               & Keys: No, Values: Yes \\
\textbf{Set}        & \code{\{1, 2, 3\}}    & No                & Yes               & No (Unique only)      \\ \bottomrule
\end{tabular}
\end{table}

\paragraph{Lists in Action: Indexing and Slicing}
Lists are created with square brackets \code{[]}. Python uses \textbf{0-based indexing}, meaning the first element is at index \code{0}. Python also supports negative indexing, where index \code{-1} refers to the last element.
\begin{equation}
\text{Slice Syntax: } \code{sequence[start : stop : step]}
\end{equation}
The slice includes the element at \code{start} up to, but \textbf{excluding}, the element at \code{stop}.

\paragraph{Python Code}
\begin{lstlisting}[language=Python]
# 1. Lists: Tracking dynamic sensor buffers
temp_history = [21.5, 22.0, 22.8, 23.1, 24.0]
temp_history.append(24.5)     # Add to end
first_val = temp_history[0]    # 21.5 (0-indexed)
last_val = temp_history[-1]    # 24.5 (negative indexing)
middle_three = temp_history[1:4]  # Slicing: indexes 1, 2, 3 -> [22.0, 22.8, 23.1]

# 2. Tuples: Fixed hardware specifications (immutable)
esp32_specs = ("ESP32-WROOM-32", 240, 520)  # (Model, Clock MHz, SRAM KB)
model, clock_mhz, sram_kb = esp32_specs      # Tuple unpacking

# 3. Dictionaries: Key-Value telemetry payloads
sensor_payload = {
    "device_id": "NODE_A12",
    "location": "North Orchard",
    "temperature": 23.4,
    "battery_pct": 89
}
# Safe access with .get() avoids crashing if key does not exist
rssi = sensor_payload.get("rssi_dbm", -999)  # Default fallback -999

# 4. Sets: Extracting unique sensor hardware IDs from a noisy stream
raw_nodes_heard = ["ESP_01", "ESP_02", "ESP_01", "ESP_03", "ESP_02", "ESP_04"]
unique_nodes = set(raw_nodes_heard)

print(f"List Slicing (1 to 4): {middle_three}")
print(f"Unpacked Tuple: {model} with {sram_kb} KB RAM")
print(f"Dictionary Device ID: {sensor_payload['device_id']}, RSSI: {rssi}")
print(f"Unique Nodes Heard ({len(unique_nodes)} devices): {sorted(list(unique_nodes))}")
\end{lstlisting}

\begin{lstlisting}[style=outputstyle]
List Slicing (1 to 4): [22.0, 22.8, 23.1]
Unpacked Tuple: ESP32-WROOM-32 with 520 KB RAM
Dictionary Device ID: NODE_A12, RSSI: -999
Unique Nodes Heard (4 devices): ['ESP_01', 'ESP_02', 'ESP_03', 'ESP_04']
\end{lstlisting}

\begin{warning}[KeyError and Missing Keys]
Accessing a dictionary key that does not exist using bracket notation raises a fatal \code{KeyError}:
\begin{lstlisting}[language=Python]
print(sensor_payload["voltage"])  # Crashes with KeyError: 'voltage'
\end{lstlisting}
Always use \code{.get("voltage", default_value)} when dealing with external IoT data packets, where fields can be dropped during wireless transmission.
\end{warning}

\begin{takeaways}
\begin{itemize}
    \item Lists store sequences that grow, shrink, or update over time.
    \item Tuples protect fixed configurations from accidental overwrites.
    \item Dictionaries provide instant $O(1)$ lookups by human-readable keys.
    \item Sets automatically strip duplicate values from streams.
\end{itemize}
\end{takeaways}

% ==============================================================================
% 6. LIST COMPREHENSIONS AND GENERATORS
% ==============================================================================
\subsection{List Comprehensions and Generators}

\begin{keyidea}[List Comprehensions and Lazy Generators]
A \textbf{list comprehension} is an elegant, one-line syntax for building a new list by transforming or filtering elements from an existing collection.

A \textbf{generator} is an iterable that computes its next value on the fly (called \textbf{lazy evaluation}) rather than storing the entire dataset in RAM simultaneously.
\end{keyidea}

\paragraph{Why It Exists: The Real-World Problem}
Microcontrollers and edge gateways have tiny memories. An ESP32 has roughly $320\text{ KB}$ of usable RAM, while a Raspberry Pi Pico has only $264\text{ KB}$. If an environmental monitor collects $500{,}000$ accelerometer readings, loading all $500{,}000$ float values into a standard Python list takes several megabytes of RAM, crashing the device with an \code{OutOfMemoryError}. Generators allow us to stream and process millions of data points using almost zero memory.

\paragraph{The Idea in Words, Step by Step}
\textbf{List Comprehension Syntax:}
\begin{equation}
\code{[expression for item in iterable if condition]}
\end{equation}
\begin{enumerate}
    \item Starts and ends with square brackets \code{[]}.
    \item Reads items sequentially from an iterable.
    \item Evaluates the optional \code{if} filter.
    \item Evaluates \code{expression} and stores the result in a brand new list in memory.
\end{enumerate}

\textbf{Generators and \code{yield}:}
Instead of returning a complete list and shutting down, a generator function uses the \code{yield} keyword. When called, it pauses execution and yields a single value to the caller. When the caller asks for the next value, the generator wakes up exactly where it paused, computes the next item, and pauses again.

\paragraph{Python Code}
\begin{lstlisting}[language=Python]
import sys

# Raw temperature readings in Celsius
celsius_readings = [18.5, 20.2, 22.0, 24.5, 26.1, 28.0]

# List Comprehension: Convert valid warm readings (>= 22 C) to Kelvin
# Formula: K = C + 273.15
kelvin_warm = [c + 273.15 for c in celsius_readings if c >= 22.0]
print(f"Warm Readings in Kelvin: {kelvin_warm}")

# Generator Function: Simulating a live endless or large telemetry stream
def simulated_adc_stream(num_samples: int):
    """Generates simulated ADC readings one by one without allocating an array."""
    current = 2000
    for i in range(num_samples):
        # A simple simulated drift
        current = (current + (i % 7) - 3)
        yield current  # Pauses and returns single value

# Comparing memory consumption
sample_count = 100_000

# 1. Full List in memory
list_in_memory = [val for val in range(sample_count)]
# 2. Generator object
gen_in_memory = (val for val in range(sample_count))

print(f"\nRAM for {sample_count:,} items in List:      {sys.getsizeof(list_in_memory):,} bytes")
print(f"RAM for {sample_count:,} items in Generator: {sys.getsizeof(gen_in_memory):,} bytes")

# Consuming 3 items from our generator stream
stream = simulated_adc_stream(5)
print("\nPulling values from generator on demand:")
print(f"  First sample:  {next(stream)}")
print(f"  Second sample: {next(stream)}")
print(f"  Third sample:  {next(stream)}")
\end{lstlisting}

\begin{lstlisting}[style=outputstyle]
Warm Readings in Kelvin: [295.15, 297.65, 299.25, 301.15]

RAM for 100,000 items in List:      800,984 bytes
RAM for 100,000 items in Generator: 200 bytes

Pulling values from generator on demand:
  First sample:  1997
  Second sample: 1995
  Third sample:  1994
\end{lstlisting}

\paragraph{Line-by-Line Explanation}
\begin{itemize}
    \item Line 8: The list comprehension combines looping, filtering (\code{c >= 22.0}), and mathematical conversion (\code{c + 273.15}) into a single readable line.
    \item Line 16: The \code{yield} statement turns \code{simulated\_adc\_stream} into a generator. It produces one integer at a time when requested.
    \item Lines 24--25: A list of $100{,}000$ numbers consumes over $800\text{ KB}$ of RAM, whereas the equivalent generator expression consumes only $200\text{ bytes}$—a $4{,}000\times$ memory reduction!
    \item Lines 31--33: We use the built-in \code{next()} function to pull individual items from the generator one at a time.
\end{itemize}

\begin{takeaways}
\begin{itemize}
    \item List comprehensions replace multi-line loop-and-append boilerplate code.
    \item Generators use \code{yield} to stream items one at a time via lazy evaluation.
    \item For edge computing and large datasets, generators prevent memory crashes.
\end{itemize}
\end{takeaways}

% ==============================================================================
% 7. CLASSES AND OBJECTS
% ==============================================================================
\subsection{Classes and Objects}

\begin{keyidea}[Object-Oriented Programming (OOP)]
A \textbf{class} is a blueprint for creating custom data types. An \textbf{object} is a concrete instance created from that blueprint. A class bundles together:
\begin{itemize}
    \item \textbf{Attributes}: Data or state variables describing the object.
    \item \textbf{Methods}: Functions that live inside the object and operate on its attributes.
\end{itemize}
\end{keyidea}

\paragraph{Why It Exists: The Real-World Problem}
In machine learning pipelines, you do not simply pass floating-point numbers around loose. You manage complex components: an entire neural network layer, a dataset loader, or a physical IoT hardware node. You need a way to bundle the node's configuration (pin numbers, baud rate, node ID) together with its actions (read, calibrate, sleep). Classes give us this modular structure.

\begin{analogy}[The Architectural Blueprint]
Think of a class as an architectural blueprint for a house. The blueprint specifies that every house has walls, a front door, and an electric meter (attributes), and can be locked or unlocked (methods). The blueprint is not a house—you cannot sleep in a blueprint. An \textbf{object} is the actual physical house built on Maple Street using that blueprint. You can build twenty unique houses from one single blueprint.
\end{analogy}

\paragraph{The Idea in Words, Step by Step}
\begin{enumerate}
    \item You declare a class using the \code{class} keyword: \code{class SensorNode:}.
    \item The constructor method, named \code{\_\_init\_\_}, runs automatically every time a new object is created.
    \item The first parameter of any method in a class is always \code{self}. \code{self} refers to the specific instance of the object whose method is being executed.
    \item We attach attributes to the instance using dot notation: \code{self.node\_id = node\_id}.
    \item Methods can access and modify the object's internal state through \code{self}.
\end{enumerate}

\paragraph{Python Code}
\begin{lstlisting}[language=Python]
class SensorNode:
    """Represents a remote IoT environmental telemetry node."""
    
    def __init__(self, node_id: str, location: str, alert_threshold: float = 30.0):
        # Instance attributes
        self.node_id = node_id
        self.location = location
        self.alert_threshold = alert_threshold
        self.history = []  # Stores recent temperature measurements
    
    def record_measurement(self, temp: float) -> bool:
        """Stores a reading and returns True if the alert threshold was exceeded."""
        self.history.append(temp)
        is_alert = temp > self.alert_threshold
        return is_alert
    
    def get_average(self) -> float:
        """Calculates and returns the average of all recorded measurements."""
        if not self.history:
            return 0.0
        return sum(self.history) / len(self.history)

# Creating two distinct objects (instances) from the same class blueprint
node_greenhouse = SensorNode(node_id="ESP_GH_01", location="Greenhouse A", alert_threshold=32.0)
node_server_room = SensorNode(node_id="ESP_SR_02", location="Server Rack 4", alert_threshold=25.0)

# Simulating readings
node_greenhouse.record_measurement(28.5)
node_greenhouse.record_measurement(33.1)  # Exceeds 32.0!
node_server_room.record_measurement(21.0)
node_server_room.record_measurement(22.4)

print(f"[{node_greenhouse.node_id}] {node_greenhouse.location}:")
print(f"  Readings Count: {len(node_greenhouse.history)}")
print(f"  Average Temp:   {node_greenhouse.get_average():.2f} C")

print(f"\n[{node_server_room.node_id}] {node_server_room.location}:")
print(f"  Readings Count: {len(node_server_room.history)}")
print(f"  Average Temp:   {node_server_room.get_average():.2f} C")
\end{lstlisting}

\begin{lstlisting}[style=outputstyle]
[ESP_GH_01] Greenhouse A:
  Readings Count: 2
  Average Temp:   30.80 C

[ESP_SR_02] Server Rack 4:
  Readings Count: 2
  Average Temp:   21.70 C
\end{lstlisting}

\paragraph{Line-by-Line Explanation}
\begin{itemize}
    \item Line 4: The \code{\_\_init\_\_} method initializes each new instance. \code{self.history = []} ensures each node has its own private list of readings.
    \item Lines 11--15: The \code{record\_measurement} method appends data to the object's internal list and checks whether the reading exceeds that node's threshold.
    \item Lines 23--24: We instantiate two independent objects. \code{node\_greenhouse} and \code{node\_server\_room} maintain completely separate variables in memory.
\end{itemize}

\begin{takeaways}
\begin{itemize}
    \item Classes bundle state (attributes) and behavior (methods) into clean, reusable structures.
    \item The constructor \code{\_\_init\_\_} initializes instance state upon creation.
    \item \code{self} represents the current instance and must be the first parameter of any instance method.
\end{itemize}
\end{takeaways}

% ==============================================================================
% 8. FILE I/O: TEXT, JSON, AND CSV
% ==============================================================================
\subsection{File Input/Output: Text, JSON, and CSV}

\begin{keyidea}[Persistent Storage via File I/O]
\textbf{File Input/Output (I/O)} allows a program to save data permanently onto storage (hard drives, SD cards, or flash memory) and read it back later.
\begin{itemize}
    \item \textbf{Text Files (\code{.txt})}: Raw strings and logs.
    \item \textbf{CSV (\code{.csv})}: Comma-Separated Values, the standard format for tabular datasets.
    \item \textbf{JSON (\code{.json})}: JavaScript Object Notation, the standard lightweight format for structured metadata, configs, and web API responses.
\end{itemize}
\end{keyidea}

\paragraph{Why It Exists: The Real-World Problem}
When your Python program stops or an edge device resets, everything stored in RAM disappears instantly. To train models on historical sensor logs, save model weights, or export predictions, we must write data to non-volatile storage.

\paragraph{The Idea in Words, Step by Step}
Always access files using Python's \textbf{context manager}: the \code{with open(...) as f:} pattern.
\begin{enumerate}
    \item The \code{open(filename, mode)} function specifies the target file and operation mode:
    \begin{itemize}
        \item \code{"r"}: Read mode (default; fails if the file does not exist).
        \item \code{"w"}: Write mode (creates a new file or overwrites an existing file completely).
        \item \code{"a"}: Append mode (adds new data to the end without deleting existing content).
    \end{itemize}
    \item The \code{with} statement guarantees that the file is safely closed and flushed to disk immediately upon exiting the block, even if an unexpected error occurs.
\end{enumerate}

\paragraph{Python Code}
\begin{lstlisting}[language=Python]
import csv
import json

# 1. Writing and Reading Tabular Sensor Data with CSV
csv_filename = "sensor_logs.csv"
telemetry_data = [
    ["timestamp_s", "sensor_id", "temp_c", "humidity_pct"],
    [1710000000, "NODE_01", 22.4, 60.1],
    [1710000010, "NODE_01", 22.7, 59.8],
    [1710000020, "NODE_02", 24.1, 55.0]
]

# Write CSV
with open(csv_filename, mode="w", newline="", encoding="utf-8") as f:
    writer = csv.writer(f)
    writer.writerows(telemetry_data)

# Read CSV and calculate average temperature
total_temp = 0.0
count = 0
with open(csv_filename, mode="r", encoding="utf-8") as f:
    reader = csv.DictReader(f)  # Automatically parses header into dictionary keys
    for row in reader:
        total_temp += float(row["temp_c"])
        count += 1

print(f"Read {count} records from CSV. Average Temp: {total_temp / count:.2f} C")

# 2. Storing and Loading Structured Config with JSON
config_filename = "device_config.json"
device_config = {
    "network": {
        "ssid": "AgriNet_Field_5G",
        "gateway_ip": "192.168.1.1"
    },
    "sampling_interval_sec": 10,
    "calibration_offset": -0.45
}

# Write JSON
with open(config_filename, mode="w", encoding="utf-8") as f:
    json.dump(device_config, f, indent=4)

# Read JSON back into a Python dictionary
with open(config_filename, mode="r", encoding="utf-8") as f:
    loaded_config = json.load(f)

print(f"Loaded Config SSID: {loaded_config['network']['ssid']}")
print(f"Calibration Offset: {loaded_config['calibration_offset']}")
\end{lstlisting}

\begin{lstlisting}[style=outputstyle]
Read 3 records from CSV. Average Temp: 23.07 C
Loaded Config SSID: AgriNet_Field_5G
Calibration Offset: -0.45
\end{lstlisting}

\paragraph{Line-by-Line Explanation}
\begin{itemize}
    \item Line 14: \code{with open(...) as f:} cleanly opens the file. When line 16 finishes, the file handle closes automatically.
    \item Line 22: \code{csv.DictReader} reads the first line as dictionary keys, making each row accessible by column name (e.g., \code{row["temp\_c"]}).
    \item Line 38: \code{json.dump(..., indent=4)} serializes the nested dictionary into structured text with four-space indentation.
    \item Line 42: \code{json.load(f)} reads the file and reconstructs the Python dictionary with nested dictionaries and numeric types preserved.
\end{itemize}

\begin{iotbox}[Logging to Flash on ESP32 (LittleFS / SD Cards)]
On edge microcontrollers running MicroPython, you log sensor measurements directly to onboard flash memory using LittleFS or an external SPI micro-SD card module using the exact same \code{with open("/sd/log.csv", "a") as f:} syntax. Appending lines every $10\text{ seconds}$ guarantees zero data loss during cellular network outages.
\end{iotbox}

\begin{takeaways}
\begin{itemize}
    \item Always use \code{with open(...) as f:} to guarantee files close safely.
    \item CSV is best for tabular numerical streams and datasets.
    \item JSON is ideal for configuration files, hyperparameters, and structured metadata.
\end{itemize}
\end{takeaways}

% ==============================================================================
% 9. ERROR HANDLING: EXCEPTIONS AND RESILIENCE
% ==============================================================================
\subsection{Error Handling: Exceptions and Resilience}

\begin{keyidea}[Defensive Programming with Exceptions]
An \textbf{exception} is an error signal raised when something goes wrong during program execution. \textbf{Error handling} wraps fragile operations in a \code{try...except} block, allowing the program to intercept and handle the failure gracefully instead of crashing.
\end{keyidea}

\paragraph{Why It Exists: The Real-World Problem}
In the physical world, sensors fail. A loose wire drops an $\text{I}^2\text{C}$ connection, an SD card runs out of space, or a corrupted wireless packet delivers gibberish instead of a floating-point number. If an unhandled error crashes a remote solar-powered weather station or industrial monitor, someone may have to climb a radio tower or drive four hours into a forest just to press a physical reset button.

\begin{analogy}[The Electrical Circuit Breaker]
Think of error handling as the circuit breaker box in your home. If a short circuit occurs in a bedroom toaster, the breaker trips instantly. The electricity cuts out safely in that single room, preventing your entire house from catching fire. Once you unplug the toaster, you flip the switch back on. A \code{try...except} block isolates risky code so a minor failure does not bring down the entire system.
\end{analogy}

\paragraph{The Idea in Words, Step by Step}
\begin{enumerate}
    \item \code{try}: Code that might fail (parsing user input, reading hardware, opening files) is placed in this indented block.
    \item \code{except SpecificError as e}: If the named error occurs, execution jumps immediately into this block to log the issue or take corrective action.
    \item \code{else}: Runs only if the \code{try} block succeeded without raising any errors.
    \item \code{finally}: Guaranteed to run under all circumstances, whether an error occurred or not. Ideal for closing hardware pins or network sockets.
\end{enumerate}

\paragraph{Python Code}
\begin{lstlisting}[language=Python]
def parse_telemetry_packet(raw_payload: str) -> dict:
    """
    Safely parses an incoming telemetry string into clean metrics.
    Expected format: 'device_id,temperature,voltage'
    """
    try:
        parts = raw_payload.split(",")
        if len(parts) != 3:
            raise ValueError(f"Expected 3 comma-separated tokens, got {len(parts)}")
        
        node_id = parts[0].strip()
        temperature = float(parts[1].strip())
        voltage = float(parts[2].strip())
        
        # Guard against physically impossible readings
        if not (-40.0 <= temperature <= 85.0):
            raise ValueError(f"Temperature {temperature} C out of physical sensor range")
        
        metrics = {
            "node_id": node_id,
            "temp_c": temperature,
            "voltage": voltage,
            "status": "VALID"
        }
        return metrics

    except ValueError as ve:
        # Handle formatting or numeric conversion failures
        return {"status": "ERROR", "reason": f"Value Error: {ve}"}
    except Exception as e:
        # Catch-all for any other unexpected edge case
        return {"status": "ERROR", "reason": f"Unexpected failure: {e}"}
    finally:
        # Runs every time, ideal for audit trails
        pass

# Test cases: clean, corrupted text, and missing field
packet_valid = "ESP_A, 24.3, 3.28"
packet_corrupt = "ESP_B, NaN_VAL, 3.30"
packet_truncated = "ESP_C, 25.1"

for raw in [packet_valid, packet_corrupt, packet_truncated]:
    result = parse_telemetry_packet(raw)
    print(f"Input: '{raw}' -> Result: {result}")
\end{lstlisting}

\begin{lstlisting}[style=outputstyle]
Input: 'ESP_A, 24.3, 3.28' -> Result: {'node_id': 'ESP_A', 'temp_c': 24.3, 'voltage': 3.28, 'status': 'VALID'}
Input: 'ESP_B, NaN_VAL, 3.30' -> Result: {'status': 'ERROR', 'reason': "Value Error: could not convert string to float: 'NaN_VAL'"}
Input: 'ESP_C, 25.1' -> Result: {'status': 'ERROR', 'reason': 'Value Error: Expected 3 comma-separated tokens, got 2'}
\end{lstlisting}

\paragraph{Line-by-Line Explanation}
\begin{itemize}
    \item Lines 9--10: If a packet arrives with missing fields, we explicitly trigger our own exception using \code{raise ValueError(...)}.
    \item Line 13: If \code{parts[1]} contains non-numeric text like \code{"NaN\_VAL"}, \code{float()} raises a \code{ValueError}.
    \item Lines 24--26: Instead of crashing the entire script, Python jumps to the \code{except ValueError} block, bundles the failure message, and allows the loop to continue to the next packet.
\end{itemize}

\begin{warning}[Never Use a Bare \code{except:}]
Avoid writing a bare \code{except:} without an error type:
\begin{lstlisting}[language=Python]
try:
    do_something()
except:  # BAD PRACTICE!
    pass
\end{lstlisting}
This catches and suppresses \textbf{everything}, including \code{KeyboardInterrupt} (pressing \texttt{Ctrl+C} to stop your script) and syntax typos, making bugs nearly impossible to track down. Always catch specific exceptions: \code{except (ValueError, KeyError, OSError) as e:}.
\end{warning}

\begin{takeaways}
\begin{itemize}
    \item Real-world inputs and network streams are dirty; always program defensively.
    \item Use \code{try...except} to catch anticipated failures and keep long-running processes alive.
    \item Always specify the exception type you want to catch; never use bare \code{except:}.
\end{itemize}
\end{takeaways}

% ==============================================================================
% 10. CHECK YOURSELF
% ==============================================================================
\subsection{Check Yourself}

Test your understanding of the concepts covered in this section with these questions:

\begin{enumerate}
    \item \textbf{Data Types \& Mutability}: What is the primary difference between a list and a tuple in Python? Why might you prefer a tuple over a list when defining fixed hardware pin numbers for an ESP32 microcontroller?
    
    \item \textbf{Hand Calculation (Loop Iterations)}: Consider the following Python code snippet:
\begin{lstlisting}[language=Python]
accum = 10.0
readings = [2.0, 4.0, 6.0]
for x in readings:
    if x == 4.0:
        continue
    accum += x * 2.0
\end{lstlisting}
    Calculate the final value of \code{accum} step by step by hand.
    
    \item \textbf{Code Diagnosis}: The following function was written to compute the average temperature from a dictionary of sensor nodes, but it contains two distinct bugs. Identify both bugs:
\begin{lstlisting}[language=Python]
def calculate_network_mean(node_dict):
    total = 0
    for node in node_dict:
        total = total + node_dict[node]
    mean = total / len(node_dict)
    # End of function
\end{lstlisting}
    
    \item \textbf{Memory \& Generators}: Explain why running \code{sum([x for x in range(10_000_000)])} might cause a resource-constrained edge computer to crash, whereas \code{sum(x for x in range(10_000_000))} runs smoothly.
    
    \item \textbf{Defensive Logic}: You are reading raw lines from a serial port connected to a soil sensor. The line should contain a float representing moisture percentage (e.g., \texttt{"42.5"}). Write a small Python function \code{safe\_read(line)} using \code{try...except} that returns the float value if valid, or \code{0.0} if the string is corrupted or empty.
\end{enumerate}

% ==============================================================================
% ANSWERS
% ==============================================================================
\subsection*{Answers}

\begin{enumerate}
    \item \textbf{Data Types \& Mutability}:
    \begin{itemize}
        \item \textbf{Difference}: A list is \textit{mutable} (its elements can be added, removed, or changed in place after creation). A tuple is \textit{immutable} (its contents are permanent and cannot be modified after instantiation).
        \item \textbf{Hardware Pin Choice}: Microcontroller hardware pins (such as \code{I2C\_PINS = (21, 22)}) represent fixed physical circuit connections. Using an immutable tuple ensures that no function or third-party library in your software can accidentally overwrite or reassign your hardware pin mapping at runtime, preventing electrical errors.
    \end{itemize}

    \item \textbf{Hand Calculation}:
    \begin{itemize}
        \item Initial state: $\code{accum} = 10.0$.
        \item \textbf{Iteration 1} ($x = 2.0$): Condition $x == 4.0$ is \code{False}. We execute $\code{accum} = 10.0 + (2.0 \times 2.0) = 10.0 + 4.0 = 14.0$.
        \item \textbf{Iteration 2} ($x = 4.0$): Condition $x == 4.0$ is \code{True}. The \code{continue} statement executes immediately. Lines below it are skipped. \code{accum} remains $14.0$.
        \item \textbf{Iteration 3} ($x = 6.0$): Condition $x == 4.0$ is \code{False}. We execute $\code{accum} = 14.0 + (6.0 \times 2.0) = 14.0 + 12.0 = 26.0$.
        \item \textbf{Final Result}: $\code{accum} = 26.0$.
    \end{itemize}

    \item \textbf{Code Diagnosis}:
    \begin{itemize}
        \item \textbf{Bug 1 (Missing Return)}: The function calculates \code{mean}, but never returns it using the \code{return} statement. Calling this function will always evaluate to \code{None}.
        \item \textbf{Bug 2 (Division by Zero)}: If an empty dictionary (\code{\{\}}) is passed into \code{node\_dict}, \code{len(node\_dict)} evaluates to $0$. Dividing by zero raises an unhandled \code{ZeroDivisionError}, crashing the script.
    \end{itemize}

    \item \textbf{Memory \& Generators}:
    \begin{itemize}
        \item In the first expression, square brackets \code{[...]} create a full \textbf{list comprehension}. Python must allocate contiguous memory in RAM to hold all $10{,}000{,}000$ integer objects at the same time (consuming roughly $80\text{ MB}$ to $100\text{ MB}$ of memory). On memory-limited microcontrollers or edge gateways, this causes an \code{OutOfMemoryError}.
        \item In the second expression, parentheses \code{(...)} create a \textbf{generator expression}. It produces integers one at a time via lazy evaluation. The \code{sum()} function consumes each integer and discards it immediately, requiring only a few bytes of RAM regardless of whether the range is ten or ten billion.
    \end{itemize}

    \item \textbf{Defensive Logic Implementation}:
\begin{lstlisting}[language=Python]
def safe_read(line: str) -> float:
    """Safely converts a serial line to float, returning 0.0 on error."""
    try:
        return float(line.strip())
    except (ValueError, TypeError, AttributeError):
        return 0.0
\end{lstlisting}
    This implementation strips leading/trailing whitespace and newlines, attempts the conversion to \code{float}, and catches \code{ValueError} (invalid text), \code{TypeError} (wrong data type passed), and \code{AttributeError} (if \code{None} was passed), safely falling back to \code{0.0}.
\end{enumerate}
